\documentclass[12pt]{article}
\usepackage[utf8]{inputenc}
\usepackage[margin=0.75in]{geometry}
\usepackage{amsmath}
\usepackage{amssymb}
\usepackage{amsthm}
\usepackage{graphicx} % Required for inserting images
\usepackage{array}
\usepackage{tikz-cd}


\title{MAT1100 Homework 4}
\author{Aaron Avram}
\date{November 17 2026}

\begin{document}
\newtheorem{definition}{Definition}
\newtheorem{theorem}{Theorem}
\newtheorem{claim}{Claim}
\newtheorem{corollary}{Corollary}

\maketitle

\begin{align*}
    \textbf{Question 5}
\end{align*}

\begin{enumerate}
    \item[(i)]
    \begin{proof}
        To show that $p(x) = x^4 + 1$ is irreducible in $\mathbb Q[x]$
        it suffices to show that $p(x+1)$ is. To see this we can expand
        \[ p(x+1) = (x + 1)^4 + 1 = x^4 + 4x^3 + 6x^2 + 4x + 2 \]
        so we can apply Eisenstein's criterion over $\mathbb Q[x]$ at the value 2 to find that $p(x+1)$
        is irreducible and hence $p(x)$ is.
    \end{proof}

    \item[(ii)]
    \begin{proof}
        First we show that either $-2$, $-1$ or $2$ is a quadratic residue modulo $p$
        for any prime $p$. A number $a$ is a quadratic residue modulo $p$ if there exists $x$
        such that $x^2 = a$ modulo p.
        
        First if $p = 2$ then $-1 = 1$ so clearly $-1$ is a quadratic residue modulo 2.

        If $p$ is an odd prime, note that $-2, -1, 2$ are all not 0 mod $p$. Now consider the square map $\varphi: (\mathbb F_p)^\times \to (\mathbb F_p)^\times$
        given by $\varphi(x) = x^2$. Clearly $\ker \varphi = \{ \pm 1 \}$ as $x^2 = 1$ has solutions $\pm 1$
        modulo $p$ as $\mathbb F_p$ is a field and for $p > 2$ these are distinct solutions.

        Thus $|\operatorname{im} \varphi| = \frac{p-1}{2}$. Since $\mathbb F_p$ is a field
        its multiplicative group is abelian and therefore $\operatorname{im}\varphi \lhd (\mathbb F_p)^\times$.
        Now
        \[ \frac{(\mathbb F_p)^\times}{\operatorname{im}\varphi} \cong \{ \pm 1 \} \]
        where the group on the right side has multiplication as its operation. Now,
        quadratic residue's modulo p are exactly the elements whose image in the quotient is $1$.
        Moreover, since $(-1)(-1) = 1$ this implies that if $a$ and $b$ are not quadratic residues
        modulo $p$ then $ab$ is.

        Now if $-1$ or $2$ are quadratic residues modulo $p$ we are done, else
        $-1 \cdot 2 = -2$ is a quadratic residue modulo $p$.

        Now to show that $x^4 + 1$ is reducible. Case 1, $\exists a \in \mathbb F_p$
        such that $a^2 = -1$ mod p then
        \[ (x^2 + a)(x^2 + a^{-1}) = x^4 + (a + a^{-1})x^2 + 1 = x^4 + a^{-1}(a^2 + 1)x^2 + 1 = x^4 + 1 \]
        and both factors are monic of positive degree so this is a non-trivial factorization of $x^4 + 1$.
        Otherwise, if $\exists a \in \mathbb F_p$ such that $a^2 = -2$ mod p, then
        \[ (x^2 + ax - 1)(x^2 - ax -1) = x^4 + (a - a)x^3 + (-1 - 1 - a^2)x^2 + (-a + a)x + 1 = x^4 + 1 \]
        and clearly this is a non-trivial factorization.
        Finally, if $\exists a \in \mathbb F_p$ such that $a^2 = 2$ mod p then
        \[ (x^2 + ax + 1)(x^2 - ax + 1) = x^4 + (a - a)x^3 + (1 + 1 - a^2)x^2 + (a - a)x + 1 = x^4 + 1 \]
        and this is clearly a non-trivial factorization.

        Thus $x^4 + 1$ is reducible mod p for any $p$.
    \end{proof}

    \item[(iii)]
    \begin{proof}
        $\mathbb Q[x]/(x^4 + 1)$ is a field and therefore also an integral domain as $x^4 + 1$ is an irreducible polynomial
        so $(x^4 + 1)$ is a maximal ideal as if $I$ is a proper ideal of $\mathbb Q[x]$
        containing $(x^4 + 1)$ then since $\mathbb Q[x]$ is a PID it is generated by a polynomial $f(x)$ which is not a unit.
        Then $f(x) \mid x^4 + 1$ but since $f(x)$ is not a unit and $x^4 + 1$ is irreducible,
        this means $f(x)$ and $x^4 + 1$ must differ by a unit. Thus $I = (x^4 + 1)$
        and therefore $(x^4 + 1)$ is maximal and therefore the quotient by it is a field.

        For all primes $p$, $\mathbb F_p[x]/(x^4 + 1)$ is not a domain and therefore also not a field. This is clear as
        for all primes $p$ $x^4 + 1$ splits as the product of two monic quadratic polynomials
        $f$ and $g$. Clearly $f, g \neq 0$ mod $x^4 + 1$ as this would imply that
        $x^4 + 1$ divides them which is not possible as $\deg f, g = 2 < \deg (x^4 + 1) = 4$.
        However, $fg = x^4 + 1 = 0 \mod (x^4 + 1)$ so $\mathbb F_p[x]/(x^4 + 1)$
        has zero divisors and therefore is not a domain.
    \end{proof}
\end{enumerate}
\end{document}